\chapter[Group Actions and Structure]{Group Actions and Structural Group Theory}

\section{Actions, Orbits, Stabilizers, and Permutation Representations}

An action studies a group by observing the objects it moves.  This makes an
abstract group concrete without choosing generators or a multiplication table:
an action is a systematic way to realize its elements as permutations.

\begin{definition}
An \emph{action} of \(G\) on \(X\) is a map \(G\times X\to X\), written
\((g,x)\mapsto g\cdot x\), satisfying the following two requirements:
\begin{enumerate}
\setlength{\itemsep}{5pt}
\item the group identity does nothing: \(1\cdot x=x\) for all \(x\in X\);
\item successive motions combine according to multiplication:
      \((gh)\cdot x=g\cdot(h\cdot x)\).
\end{enumerate}
The \emph{orbit} of \(x\) is \(Gx=\{g\cdot x:g\in G\}\), and its
\emph{stabilizer} is \(G_x=\{g\in G:g\cdot x=x\}\).
\end{definition}

Equivalently, an action is a homomorphism \(G\to\operatorname{Sym}(X)\).
This equivalence explains the two action axioms: they say exactly that the
identity acts as the identity permutation and a product acts by composition.
The natural action of \(S_n\) on \(\{1,\ldots,n\}\), the action of \(D_n\)
on the vertices of a regular polygon, and the left-translation action of \(G\)
on itself are the basic examples.

It is useful to pause over these examples.  The natural action of \(S_n\) is
transitive, and the stabilizer of \(1\) consists of permutations of the other
\(n-1\) letters.  The geometric action of \(D_n\) is also transitive; the
stabilizer of a vertex has two elements, the identity and the reflection through
that vertex.  The left regular action of \(G\) is transitive with trivial
stabilizer.  Finally, the coset action on \(G/H\) is transitive and has
stabilizer \(H\) at the coset \(H\).  These are the examples to test whenever
an abstract action appears later.

\begin{definition}
Let \(\rho:G\to\operatorname{Sym}(X)\) be the homomorphism associated to an
action.  Its \emph{kernel} is
\[\ker\rho=\{g\in G:g\cdot x=x\text{ for every }x\in X\}.\]
The action is \emph{faithful} if this kernel is trivial.
\end{definition}

Thus an action is faithful exactly when no nonidentity group element is
invisible on \(X\), equivalently when \(\rho\) is injective.  The natural
action of \(S_n\) and the left regular action of any group are faithful.  By
contrast, the action of \(D_4\) on the set of diagonals of a square is not
faithful: distinct geometric symmetries can induce the same diagonal action.

\begin{proposition}
The stabilizer is a subgroup, the orbits partition \(X\), and an action gives
a homomorphism \(G\to\operatorname{Sym}(X)\).
\end{proposition}
\begin{proof}
For the stabilizer, \(1\cdot x=x\), so \(1\in G_x\).  If \(g,h\in G_x\),
then \((gh^{-1})\cdot x=g\cdot(h^{-1}\cdot x)=g\cdot x=x\), so the subgroup
test applies.  If two orbits meet, say \(g\cdot x=h\cdot y\), then applying
\(g^{-1}\) yields \(x=g^{-1}h\cdot y\).  Hence \(x\in Gy\), and applying
group elements shows \(Gx\subseteq Gy\); symmetry gives equality.  Finally,
the assignment \(g\mapsto\rho(g)\), where \(\rho(g)(x)=g\cdot x\), satisfies
\(\rho(gh)=\rho(g)\rho(h)\) by the second action axiom.
\end{proof}

\begin{theorem}[Cayley]
Every group \(G\) is isomorphic to a subgroup of \(\operatorname{Sym}(G)\).
\end{theorem}
\begin{proof}
Let \(G\) act on itself by left translation, \(g\cdot x=gx\).  The associated
homomorphism sends \(g\) to the permutation \(x\mapsto gx\).  If this
permutation is the identity, then evaluating it at \(1\) gives \(g=1\).
Its kernel is therefore trivial, so the homomorphism is injective.
\end{proof}

\begin{example}
For the left action of \(G\) on itself, the orbit of \(1\) is all of \(G\)
and its stabilizer is trivial.  For the conjugation action, the orbit and
stabilizer become a conjugacy class and a centralizer.  These examples show
that one definition organizes both the regular action and the internal
symmetry of a group.
\end{example}

\begin{example}[The Coset Action]
For \(H\leq G\), left multiplication defines \(g\cdot(aH)=(ga)H\) on
\(G/H\).  The action is transitive and the stabilizer of \(H\) is \(H\).
Its kernel is
\[\bigcap_{g\in G}gHg^{-1}.\]
Indeed, an element fixes every coset exactly when it belongs to every conjugate
of \(H\).  This intersection, called the \emph{core} of \(H\), is the largest
normal subgroup of \(G\) contained in \(H\).  The example is an important way
to construct permutation representations from subgroups.
\end{example}

\begin{theorem}[Orbit--Stabilizer]
For a finite group action, \(|Gx|=[G:G_x]\), and therefore
\(|G|=|Gx|\,|G_x|\).
\end{theorem}
\begin{proofidea}
The cosets of \(G_x\) are exactly the different ways to move \(x\).  Two
elements \(g,h\) send \(x\) to the same point precisely when
\(h^{-1}g\) fixes \(x\), which is the equivalence relation defining the cosets.
\end{proofidea}
\begin{proof}
Define \(\Phi:G/G_x\to Gx\) by \(\Phi(gG_x)=g\cdot x\).  To check
well-definedness, suppose \(gG_x=hG_x\).  Then \(h^{-1}g\in G_x\), whence
\(g\cdot x=h\cdot((h^{-1}g)\cdot x)=h\cdot x\).  Conversely, if
\(g\cdot x=h\cdot x\), applying \(h^{-1}\) gives
\((h^{-1}g)\cdot x=x\); thus \(h^{-1}g\in G_x\) and \(gG_x=hG_x\).  This
proves injectivity at the same time.  Surjectivity holds because every point
of the orbit has the form \(g\cdot x\).  Therefore \(|Gx|=[G:G_x]\), and
Lagrange's theorem yields \(|G|=|Gx|\,|G_x|\).
\end{proof}

\begin{example}
Let \(D_4\) act on the four vertices of a square.  The orbit of any vertex is
the whole vertex set, so it has size four.  Its stabilizer contains the identity
and the reflection across the diagonal through that vertex, so it has size two.
Thus \(|D_4|=4\cdot2\), as orbit--stabilizer predicts.
\end{example}

\begin{corollary}[Burnside's Counting Lemma]
The number of orbits of a finite action on \(X\) is
\[\frac{1}{|G|}\sum_{g\in G}|\{x\in X:g\cdot x=x\}|.\]
\end{corollary}
\begin{proof}
Count pairs \((g,x)\) with \(g\cdot x=x\).  Counting by \(g\) gives the
numerator.  Counting by orbits gives \(|G|\) for each orbit, by
orbit--stabilizer.
\end{proof}

\begin{example}[Counting Colourings up to Symmetry]
Let the rotation group \(C_4\) act on the two-colourings of the vertices of a
square.  The identity fixes all \(2^4\) colourings, a quarter-turn fixes two,
the half-turn fixes four, and the other quarter-turn fixes two.  Burnside's
lemma gives \((16+2+4+2)/4=6\) rotational symmetry types.  The point is not
the answer itself: the action replaces a potentially awkward case distinction
by a fixed-point count.
\end{example}

\section{Conjugation, Centralizers, and the Class Equation}

Conjugation answers the question ``which elements look the same from inside
the group?''  It preserves order and many other structural properties.

\begin{definition}
For a group \(G\) and an element \(x\in G\), the \emph{centralizer} of \(x\)
is
\[C_G(x)=\{g\in G:gx=xg\}.\]
The \emph{center} of \(G\) is
\[Z(G)=\{g\in G:gx=xg\text{ for every }x\in G\};\]
thus \(Z(G)\) is the collection of elements that commute with all of \(G\).
For a subgroup \(H\leq G\), the \emph{normalizer} of \(H\) is
\[N_G(H)=\{g\in G:gHg^{-1}=H\}.\]
It consists of the elements whose conjugation preserves \(H\) as a set.
\end{definition}

Each of these sets is a subgroup of \(G\).  For example, the subgroup test
for \(C_G(x)\) follows from the fact that if \(a,b\) commute with \(x\), then
\((ab^{-1})x=a(b^{-1}x)=a(xb^{-1})=(ax)b^{-1}=x(ab^{-1})\).
The same calculation with every \(x\) gives the assertion for \(Z(G)\), and
the identity \((ab^{-1})H(ab^{-1})^{-1}=H\) gives it for \(N_G(H)\).

Conjugation, \(g\cdot x=gxg^{-1}\), has conjugacy classes as its orbits and
centralizer \(C_G(x)\) as the stabilizer.  The elements fixed by all of \(G\)
are precisely the center \(Z(G)\).  The normalizer similarly records the
largest subgroup of \(G\) in which \(H\) is normal.

\begin{lemma}[\(p\)-Group Action Congruence]
If a finite \(p\)-group \(P\) acts on a finite set \(X\), then
\[|X|\equiv|X^P|\pmod p,\]
where \(X^P=\{x\in X:g\cdot x=x\text{ for every }g\in P\}\).
\end{lemma}
\begin{proof}
Partition \(X\) into \(P\)-orbits.  Orbit--stabilizer makes every orbit size
a power of \(p\).  The orbits of size one are exactly the points fixed by all
of \(P\); every other orbit has size divisible by \(p\).  Adding orbit sizes
gives the stated congruence.
\end{proof}

This elementary congruence is a reusable engine: it converts a group action
into the existence of fixed points whenever \(|X|\) is not divisible by \(p\).

\begin{example}[A Fixed Point Forced Modulo \(p\)]
Let a group of order \(p\) act on a set with \(p+1\) elements.  The congruence
gives \(|X^P|\equiv p+1\equiv1\pmod p\), so there is a global fixed point.
The same orbit-counting idea powers the center theorem and Sylow congruence.
\end{example}

\begin{remark}[How to Design an Action Argument]
An action proof begins by choosing a set whose points encode the objects one
wants to study.  To prove a divisibility statement, choose a set with an
accessible cardinality and compare orbit sizes.  To prove a containment
statement, let a subgroup act on a coset space, so that a fixed coset translates
into a conjugate containment.  To study elements inside one group, use
conjugation: its orbits are conjugacy classes and its stabilizers are
centralizers.  In each case the formal calculation is short only after the
right action and the right set have been chosen.
\end{remark}

\begin{theorem}[Class Equation]
If \(x_1,\ldots,x_r\) represent the noncentral conjugacy classes of finite
\(G\), then
\[
|G|=|Z(G)|+\sum_i[G:C_G(x_i)].
\]
\end{theorem}
\begin{proof}
Partition \(G\) into conjugacy classes.  Central elements yield singleton
classes, and orbit--stabilizer gives the size of every other class.
\end{proof}

\begin{example}
In \(S_3\), the three transpositions form one conjugacy class, the two
three-cycles form another, and the identity is central.  Thus
\(6=1+3+2\).  This small example already illustrates the class equation as a
partition into symmetry types rather than a numerical trick.
\end{example}

\begin{example}[Centralizers in \(S_3\)]
The preceding equality can be recovered directly from centralizers.  The
centralizer of \((12)\) is \(\{1,(12)\}\), so orbit--stabilizer gives a
conjugacy class of size \([S_3:C_{S_3}((12))]=3\).  The centralizer of
\((123)\) is \(\langle(123)\rangle\), so its class has size two.  This is a
useful practical lesson: to count conjugates of an element, compute the
subgroup of elements that commute with it.
\end{example}

\begin{corollary}
Every finite \(p\)-group has nontrivial center.
\end{corollary}
\begin{proofidea}
For a \(p\)-group, every noncentral orbit under conjugation has size divisible
by \(p\).  The class equation therefore forces the fixed-point part, the
center, to have size divisible by \(p\) as well.
\end{proofidea}
\begin{proof}
For a noncentral \(x\), the centralizer \(C_G(x)\) is proper.  Its order is a
proper divisor of the order of the \(p\)-group \(G\), hence
\([G:C_G(x)]\) is divisible by \(p\).  Therefore the class equation gives
\[|G|=|Z(G)|+\sum_i[G:C_G(x_i)]\equiv|Z(G)|\pmod p.\]
Since \(p\mid|G|\), also \(p\mid|Z(G)|\).  In particular the center has more
than its identity element.
\end{proof}

\section{The Sylow Theorems and Structural Applications}

Sylow theory turns the divisibility of \(|G|\) into actual subgroups.  Its
proofs are a model use of actions: choose a finite set on which a \(p\)-group
acts, observe that nontrivial orbit sizes are divisible by \(p\), and extract
a fixed point from a count not divisible by \(p\).

\begin{theorem}[Cauchy]
If a prime \(p\) divides the order of a finite group \(G\), then \(G\) has an
element of order \(p\).
\end{theorem}
The unusual-looking tuple set is chosen so that it has a readily computed size
and a natural cyclic symmetry.  Once the first \(p-1\) entries are selected,
the last must be \((g_1\cdots g_{p-1})^{-1}\), so there are exactly
\(|G|^{p-1}\) tuples.  Rotation preserves the condition that the product is
one because a cyclic rotation conjugates the total product by its first entry.
\begin{proofidea}
The set of tuples has cardinality divisible by \(p\), while cyclic rotation
organizes it into orbits of size \(1\) or \(p\).  The fixed points are exactly
the repeated tuples, which encode elements satisfying \(g^p=1\).
\end{proofidea}
\begin{proof}
Let \(X\) be the set of \(p\)-tuples \((g_1,\ldots,g_p)\) with
\(g_1\cdots g_p=1\).  Its first \(p-1\) entries may be chosen freely, so
\(|X|=|G|^{p-1}\), which is divisible by \(p\).  Cyclically rotate these
tuples.  Every orbit has size \(1\) or \(p\), and the fixed tuples are exactly
\((g,\ldots,g)\) with \(g^p=1\).  Thus the number of such \(g\) is divisible
by \(p\).  It includes \(1\), so there is a nonidentity one; its order divides
the prime \(p\), hence equals \(p\).
\end{proof}

This proof illustrates a recurring discovery principle: construct a finite set
whose cardinality is divisible by \(p\), then let a group of order \(p\) act so
that its fixed points encode the desired elements.

\begin{theorem}[Sylow I: existence]
Let \(|G|=p^am\), where \(p\) is prime and \(p\nmid m\).  Then \(G\) has a
subgroup of order \(p^a\).  Such a subgroup is called a \emph{Sylow
\(p\)-subgroup}.
\end{theorem}
\begin{proofidea}
Induct on the group order.  A central element of order \(p\) lets us reduce to
a quotient.  If the center does not supply one, the class equation locates a
proper centralizer that still contains the full \(p\)-part of \(|G|\).
\end{proofidea}
\begin{proof}
We use strong induction on \(|G|\); the assertion is immediate when \(a=0\).
First suppose that \(p\mid |Z(G)|\).  By Cauchy's theorem, \(Z(G)\) contains
a subgroup \(C\) of order \(p\).  It is normal in \(G\), and
\(|G/C|=p^{a-1}m<|G|\).  Let \(\pi:G\to G/C\) be the quotient map.  The
induction hypothesis applied to \(G/C\) gives a subgroup
\(\overline P\leq G/C\) of order \(p^{a-1}\).  Define
\(P=\pi^{-1}(\overline P)\).  Then \(P/C\cong\overline P\), so
\(|P|=|C|\,|\overline P|=p\cdot p^{a-1}=p^a\).

Now suppose \(p\nmid |Z(G)|\).  In the class equation, the noncentral terms
are conjugacy-class sizes \([G:C_G(x_i)]\).  Since \(|G|\equiv0\pmod p\) but
\(|Z(G)|\not\equiv0\pmod p\), at least one of these indices is not divisible
by \(p\).  For that noncentral \(x_i\), the centralizer \(C_G(x_i)\) is a
proper subgroup of \(G\), while
\[
 |C_G(x_i)|=\frac{|G|}{[G:C_G(x_i)]}
\]
still has \(p^a\) as its full \(p\)-part.  The induction hypothesis inside
this strictly smaller group produces a subgroup of order \(p^a\), which is
also a subgroup of \(G\).
\end{proof}

\begin{theorem}[Sylow II: containment and conjugacy]
Let \(Q\) be a Sylow \(p\)-subgroup of a finite group \(G\).  Every
\(p\)-subgroup \(P\leq G\) has a conjugate contained in \(Q\).  In particular,
all Sylow \(p\)-subgroups are conjugate.
\end{theorem}
\begin{proof}
Let \(P\) act by left multiplication on the coset set \(G/Q\).  Its size is
\([G:Q]=m\), which is not divisible by \(p\).  Every \(P\)-orbit has size a
power of \(p\), by orbit--stabilizer.  If there were no fixed coset, the sum
of the orbit sizes would be divisible by \(p\), contrary to \(p\nmid m\).
Thus some coset \(gQ\) is fixed by every element of \(P\).  The equality
\(pgQ=gQ\) for all \(p\in P\) is equivalent to
\(g^{-1}pg\in Q\) for all \(p\in P\), hence to \(g^{-1}Pg\leq Q\).
If \(P\) is itself Sylow, the two groups in this containment have the same
finite order, so they are equal.  This proves conjugacy.
\end{proof}

\begin{theorem}[Sylow III: number of Sylow subgroups]
If \(|G|=p^am\) with \(p\nmid m\), the number \(n_p\) of Sylow
\(p\)-subgroups divides \(m\) and satisfies \(n_p\equiv1\pmod p\).
\end{theorem}
\begin{proof}
Fix a Sylow subgroup \(Q\).  Sylow II makes conjugation by \(G\) transitive
on the set \(\mathcal S\) of Sylow \(p\)-subgroups, so orbit--stabilizer gives
\[n_p=|\mathcal S|=[G:N_G(Q)].\]
Because \(Q\leq N_G(Q)\), the index divides \(|G|/|Q|=m\).

Now let \(Q\) act by conjugation on \(\mathcal S\).  Suppose that \(R\in
\mathcal S\) is fixed.  Then \(Q\leq N_G(R)\); since \(R\trianglelefteq
N_G(R)\), the product \(QR\) is a subgroup.  It is a \(p\)-subgroup
containing the Sylow subgroup \(R\), hence \(QR=R\), and therefore \(Q\leq
R\).  Equal orders give \(Q=R\).  Thus \(Q\) is the unique fixed point of
this action.  Every remaining orbit has size divisible by \(p\), so
\(|\mathcal S|\equiv1\pmod p\).
\end{proof}

\begin{example}
If \(|G|=21\), then \(n_7\) divides \(3\) and is congruent to \(1\) modulo
\(7\); hence \(n_7=1\), so the Sylow \(7\)-subgroup is normal.  This is a
structural conclusion, not a classification of all groups of order \(21\).
\end{example}

\begin{corollary}
A Sylow \(p\)-subgroup is normal if and only if it is the unique Sylow
\(p\)-subgroup.  In particular, every group of order \(45\) has a normal
Sylow \(5\)-subgroup.
\end{corollary}
\begin{proof}
Conjugation carries a Sylow \(p\)-subgroup to another one.  Thus a normal one
is fixed by all conjugations and, by Sylow II, is the only one.  Conversely,
a unique Sylow subgroup is fixed under conjugation, so it is normal.  For
\(|G|=45=5\cdot9\), Sylow III says \(n_5\mid9\) and \(n_5\equiv1\pmod5\);
among \(1,3,9\), only \(1\) has this congruence.
\end{proof}

\section{Semidirect Products and Split Extensions}

A direct product describes two components that commute and have no interaction.
A semidirect product records the next possibility: the second component acts on
the first by automorphisms.  The homomorphism \(\alpha:H\to\operatorname{Aut}(N)\)
is therefore not auxiliary data; it is the rule describing that interaction.

The multiplication formula is forced by the internal situation.  Suppose
\(N\trianglelefteq G\), \(H\leq G\), and every element of \(G\) has a unique
form \(nh\).  Then
\[
 (nh)(n'h')=n\bigl(hn'h^{-1}\bigr)hh'.
\]
Normality puts \(hn'h^{-1}\) back in \(N\), while the element \(h\) acts on
\(N\) by conjugation.  The external definition simply preserves this rule
when the two groups have not yet been placed inside a common ambient group.

\begin{definition}
For \(\alpha:H\to\operatorname{Aut}(N)\), the \emph{semidirect product}
\(N\rtimes_\alpha H\) is \(N\times H\) with
\((n,h)(n',h')=(n\alpha(h)(n'),hh')\).
\end{definition}

\begin{proposition}
This operation makes \(N\rtimes_\alpha H\) a group; \(N\times\{1\}\) is
normal, \(\{1\}\times H\) is a subgroup, and they have trivial intersection.
\end{proposition}
\begin{proof}
The identity is \((1,1)\).  Since every \(\alpha(h)\) is an automorphism,
\[(n,h)^{-1}=(\alpha(h^{-1})(n^{-1}),h^{-1}).\]
For associativity, the first coordinates of the two possible products are
\begin{align*}
 ((n,h)(n',h'))(n'',h'')&:
 n\alpha(h)(n')\alpha(hh')(n''),\\
 (n,h)((n',h')(n'',h''))&:
 n\alpha(h)\bigl(n'\alpha(h')(n'')\bigr).
\end{align*}
They agree because \(\alpha(hh')=\alpha(h)\alpha(h')\).  The two displayed subsets have the
claimed subgroup properties.  Finally,
\[
 (n,h)(n',1)(n,h)^{-1}=
 \bigl(n\alpha(h)(n')n^{-1},1\bigr)\in N\times\{1\},
\]
so \(N\times\{1\}\) is normal.  The two subgroups meet only in \((1,1)\).
\end{proof}

An internal semidirect product is equivalently a split short exact sequence
\[1\longrightarrow N\longrightarrow G\longrightarrow H\longrightarrow1.\]
Here ``split'' means that the quotient map \(q:G\to H\) has a homomorphic
section \(s:H\to G\). Then \(q\) restricted to \(s(H)\) is inverse to \(s\),
so \(N\cap s(H)=1\), where \(N=\ker q\). Moreover, if \(g\in G\) and
\(h=q(g)\), then
\[
q\bigl(gs(h)^{-1}\bigr)=hh^{-1}=1,
\]
so \(gs(h)^{-1}\in N\) and \(g\in Ns(H)\). Hence \(G=Ns(H)\), and
Proposition~\ref{prop:internal-semiproduct} applies. Conversely, an internal
semidirect product gives \(G/N\cong H\), and the inclusion of the complement
gives a homomorphic section. This is only a first encounter with extensions,
but it explains why semidirect products are ubiquitous.

\begin{example}
Let \(C_2=\langle s\rangle\) act on \(C_3=\langle r\rangle\) by inversion,
\(srs^{-1}=r^{-1}\).  Then \(C_3\rtimes C_2\) has presentation
\(\langle r,s\mid r^3=s^2=1,\ srs^{-1}=r^{-1}\rangle\), and the map sending
\(r\) to \((123)\) and \(s\) to \((12)\) identifies it with \(S_3\).  It is
not a direct product because its two factors do not commute.
\end{example}

\begin{example}
More generally, \(D_n\cong C_n\rtimes C_2\), where the nontrivial element of
\(C_2\) acts on \(C_n\) by inversion.  The relation \(srs^{-1}=r^{-1}\) is
exactly the twisting encoded by this action.  Thus the dihedral examples
introduced in Chapter~2 provide a geometric model for every semidirect product
in which a reflection reverses a cyclic rotation.
\end{example}

\begin{proposition}[Internal Semidirect Product]\label{prop:internal-semiproduct}
If \(N\trianglelefteq G\), \(H\leq G\), \(G=NH\), and \(N\cap H=1\), then
\(G\cong N\rtimes H\), where \(H\) acts on \(N\) by conjugation.
\end{proposition}
\begin{proof}
Every element has an expression \(nh\).  It is unique: if \(nh=n'h'\), then
\(n^{-1}n'=hh'^{-1}\) belongs to both \(N\) and \(H\), hence is \(1\).
Sending \((n,h)\) to \(nh\) is therefore bijective.  Its compatibility with multiplication is
\(nhn'h'=n(hn'h^{-1})hh'\), which is exactly the external semidirect-product
law for conjugation.
\end{proof}

\section{Composition Series, Solvability, and Nilpotence}

Composition series make precise the idea of breaking a group into indivisible
layers.  A simple group has no nontrivial proper normal subgroup, so it cannot
be decomposed further by quotients.  Jordan--H\"older says that, although the
intermediate subgroups may differ, the resulting simple layers are intrinsic.

\begin{definition}
A \emph{composition series} is a finite normal series with simple factors.  A
group is \emph{solvable} when it has a finite normal series with abelian factors.
\end{definition}

\begin{example}
The chain \(S_3\trianglerighteq A_3\trianglerighteq1\) is a composition
series, with factors \(C_2\) and \(C_3\).  The order in which the factors are
listed has no intrinsic significance; their isomorphism types do.
\end{example}

\begin{proposition}
Every finite group has a composition series.
\end{proposition}
\begin{proof}
Induct on the group order.  If \(G\ne1\), choose a maximal proper normal
subgroup \(N\), which exists because the group is finite.  Then \(G/N\) is
simple, and induction gives a composition series for \(N\).  Prepending \(G\)
to that series gives one for \(G\).
\end{proof}

\begin{theorem}[Jordan--H\"older]
Any two composition series of a finite group have the same length and the same
simple factors, up to ordering and isomorphism.
\end{theorem}
\begin{proof}
We argue by induction on \(|G|\).  There is nothing to prove for \(G=1\).
Write the two series as
\[
 G=G_0\trianglerighteq G_1=N\trianglerighteq\cdots\trianglerighteq1,
 \qquad
 G=H_0\trianglerighteq H_1=M\trianglerighteq\cdots\trianglerighteq1.
\]
The quotients \(G/N\) and \(G/M\) are simple, so \(N\) and \(M\) are maximal
proper normal subgroups of \(G\).

If \(N=M\), the tails are composition series of the smaller group \(N\).
The induction hypothesis says that these tails have the same length and the
same factors, and adding their common first factor \(G/N\) proves the claim.

Suppose \(N\ne M\).  The product \(NM\) is normal in \(G\): for \(g\in G\),
\(g(NM)g^{-1}=(gNg^{-1})(gMg^{-1})=NM\).  It strictly contains \(N\), since
otherwise \(M\leq N\), and then the maximality of \(M\) would force
\(N=M\).  Maximality of \(N\) consequently gives \(NM=G\).  Put
\(I=N\cap M\).  Since both \(N\) and \(M\) are normal in \(G\), so is \(I\),
and in particular \(I\) is normal in both \(N\) and \(M\).  Applying the Second
Isomorphism Theorem in the forms
\[
 M/(M\cap N)\cong MN/N=G/N,\qquad
 N/(N\cap M)\cong NM/M=G/M
\]
shows that \(M/I\) and \(N/I\) are simple.

Choose a composition series from \(I\) down to \(1\).  Because \(N/I\) is
simple, adjoining \(N\) as a new initial term produces a composition series
of \(N\), whose new factor is \(N/I\); similarly, adjoining \(M\) produces a
composition series of \(M\) with new factor \(M/I\).  By induction, the first constructed series has the same factors as
the tail of the original \(G\)-series beginning at \(N\), and the second has
the same factors as the tail beginning at \(M\).  The complete first factor
list is therefore the factors below \(I\), followed by \(N/I\) and \(G/N\);
the second is the same factors below \(I\), followed by \(M/I\) and \(G/M\).
The displayed isomorphisms identify \(N/I\cong G/M\) and
\(M/I\cong G/N\), so the two lists agree up to interchanging these final two
factors.
\end{proof}

For Galois theory it is useful to recognize solvability without guessing a
normal series.  The commutator \([x,y]=xyx^{-1}y^{-1}\) measures the failure
of \(x\) and \(y\) to commute: \([x,y]=1\) if and only if \(xy=yx\).  Indeed,
\(xyx^{-1}y^{-1}=1\) is equivalent, after multiplying on the right by \(yx\),
to \(xy=yx\).  The subgroup generated by all commutators is the
\emph{derived subgroup} \([G,G]\); define \(G^{(0)}=G\) and
\(G^{(i+1)}=[G^{(i)},G^{(i)}]\).

For later quotient constructions, it matters that \([G,G]\) is normal.
Conjugation sends a commutator to a commutator:
\[
 g[x,y]g^{-1}=[gxg^{-1},gyg^{-1}].
\]
Thus conjugation preserves the subgroup generated by all commutators, and
\([G,G]\trianglelefteq G\).

\begin{proposition}
The quotient \(G/[G,G]\) is abelian and is the largest abelian quotient of
\(G\): every homomorphism from \(G\) to an abelian group factors uniquely
through it.
\end{proposition}
\begin{proof}
Every commutator becomes trivial in the quotient, so its cosets commute.  If
\(f:G\to A\) and \(A\) is abelian, then \(f([x,y])=1\) for all \(x,y\), hence
\([G,G]\leq\ker f\).  The universal property of the quotient map now gives
the unique factorization.
\end{proof}

\begin{example}[The Derived Subgroup of \(S_3\)]
The quotient \(S_3/A_3\cong C_2\) is abelian, so the universal property above
gives \([S_3,S_3]\leq A_3\).  On the other hand, direct calculation gives
\[
 [(12),(23)]=(12)(23)(12)(23)=(123)^2=(132).
\]
This nonidentity element generates \(A_3\), so \(A_3\leq[S_3,S_3]\).  Hence
\([S_3,S_3]=A_3\), and \(S_3/[S_3,S_3]\cong C_2\).
\end{example}

\begin{theorem}
A group \(G\) is solvable if and only if \(G^{(r)}=1\) for some \(r\geq0\).
\end{theorem}
\begin{proof}
If the derived series reaches \(1\), its successive quotients are abelian by
the preceding proposition.  Conversely, let
\(G=G_0\trianglerighteq\cdots\trianglerighteq G_r=1\) have abelian factors.
Inductively, \(G^{(i)}\leq G_i\): if \(G^{(i)}\leq G_i\), then the image of
\(G^{(i)}\) in the abelian group \(G_i/G_{i+1}\) has trivial derived subgroup,
so \(G^{(i+1)}\leq G_{i+1}\).  Hence \(G^{(r)}=1\).
\end{proof}

Solvability asks for a group built by successive abelian quotients.  Nilpotence
is stronger: it asks that the construction proceed from the center outward.
This distinction matters already for \(S_3\), which is solvable but not
nilpotent.

\begin{definition}
Set \(Z_0(G)=1\) and define the \emph{upper central series} inductively by
\[Z_{i+1}(G)/Z_i(G)=Z\bigl(G/Z_i(G)\bigr).\]
The group \(G\) is \emph{nilpotent} if \(Z_c(G)=G\) for some \(c\).  The least
such \(c\) is its nilpotence class (unless \(G=1\)).
\end{definition}

Equivalently, a nilpotent group admits a finite chain
\[
1=N_0\leq N_1\leq\cdots\leq N_c=G
\]
with \(N_{i+1}/N_i\leq Z(G/N_i)\).  Indeed, inductively
\(N_i\leq Z_i(G)\): this is clear for \(i=0\), and centrality of
\(N_{i+1}/N_i\) gives \(N_{i+1}\leq Z_{i+1}(G)\).  Thus
\(G=N_c\leq Z_c(G)\), so \(G\) is nilpotent.  Conversely, the upper central
series itself is such a chain.  The upper central series is therefore the
canonical maximal central series, obtained by repeatedly pulling back the
center of the current quotient.  Thus abelian groups are nilpotent of class
at most one.  A nontrivial nilpotent group has nontrivial center; therefore
\(S_3\), whose center is trivial, is not nilpotent.

\begin{theorem}
Every finite \(p\)-group is nilpotent.
\end{theorem}
\begin{proof}
Induct on \(|G|\).  The center of a nontrivial finite \(p\)-group is
nontrivial, and Cauchy's theorem gives a central subgroup \(C\) of order
\(p\).  By induction, \(G/C\) has an upper central series
\[
 1=\overline Z_0\leq\overline Z_1\leq\cdots\leq\overline Z_c=G/C.
\]
Let \(Z_i\) be the inverse image of \(\overline Z_i\) in \(G\).  Then
\(Z_0=C\), and the chain
\[
 1\leq C=Z_0\leq Z_1\leq\cdots\leq Z_c=G
\]
is central: if \(z\in Z_{i+1}\) and \(g\in G\), then the image of \(z\)
is central modulo \(\overline Z_i\), so \([z,g]\in Z_i\).  The first step
is central because \(C\leq Z(G)\).  Thus this is a central series from
\(1\) to \(G\), proving that \(G\) is nilpotent.
\end{proof}

\begin{proposition}
Every nilpotent group is solvable.
\end{proposition}
\begin{proof}
In the upper central series, each factor
\(Z_{i+1}(G)/Z_i(G)\) lies in the center of \(G/Z_i(G)\), so it is abelian.
Thus the upper central series is a solvable series.
\end{proof}

\begin{proposition}[Subgroups and quotients]
Subgroups and quotients of solvable groups are solvable.
\end{proposition}
\begin{proof}
Let \(G=G_0\trianglerighteq G_1\trianglerighteq\cdots\trianglerighteq
G_r=1\) be a solvable series, and let \(H\leq G\).  Intersecting gives
\[
 H=H\cap G_0\trianglerighteq H\cap G_1\trianglerighteq\cdots
 \trianglerighteq H\cap G_r=1.
\]
Each term is normal in the preceding one.  The map
\(H\cap G_i\to G_i/G_{i+1}\) has kernel \(H\cap G_{i+1}\), so its quotient
\((H\cap G_i)/(H\cap G_{i+1})\) is isomorphic to a subgroup of the abelian
group \(G_i/G_{i+1}\).  Removing repeated terms leaves a solvable series for
\(H\).

If \(q:G\to Q\) is onto, the chain
\(Q=q(G_0)\trianglerighteq q(G_1)\trianglerighteq\cdots\trianglerighteq
q(G_r)=1\) has normal terms.  Each nontrivial factor is an image of a
subgroup of some abelian factor \(G_i/G_{i+1}\), and is therefore abelian.
After repetitions are removed, this is a solvable series for \(Q\).
\end{proof}

\begin{proposition}[Extensions]
If \(N\trianglelefteq G\), and both \(N\) and \(G/N\) are solvable, then
\(G\) is solvable.
\end{proposition}
\begin{proof}
Take a solvable series for \(N\).  Take also a solvable series for \(G/N\)
and pull its terms back under the quotient map \(G\to G/N\).  Splicing the
first series to these inverse images gives a normal series for \(G\); its
factors are, respectively, factors from the series for \(N\) and factors
isomorphic to those from the series for \(G/N\).  They are all abelian.
\end{proof}

\begin{remark}[Looking Ahead]
The action of a Galois group on roots and intermediate fields will make these
methods central again.  Solvability of groups eventually characterizes
solvability of polynomial equations by radicals.
\end{remark}

\section{Two Recurring Structural Laboratories}

The abstract constructions of this chapter become much easier to retain when
they are repeatedly tested on a few familiar groups.  The following two
examples summarize how the same group can illuminate several different ideas.

\subsection*{The Group \(S_3\)}

The six-element symmetric group is the smallest nonabelian group.  It acts
faithfully and transitively on three letters.  Orbit--stabilizer at the letter
\(1\) gives
\[|S_3|=|S_3\cdot1|\,|(S_3)_1|=3\cdot2.\]
The stabilizer is \(\{1,(23)\}\), a copy of \(C_2\).  Thus the action turns a
concrete observation about permutations into a coset calculation.

Under conjugation, its elements separate into the classes
\[\{1\},\qquad\{(12),(13),(23)\},\qquad\{(123),(132)\}.
\]
The class equation is \(6=1+3+2\).  This is more than bookkeeping: it shows
that the three transpositions are structurally indistinguishable within
\(S_3\), whereas a transposition subgroup is not normal because conjugation
moves it to another such subgroup.

The subgroup \(A_3=\langle(123)\rangle\) is normal, and the quotient
\(S_3/A_3\cong C_2\) records parity.  The subgroup \(A_3\) is complemented by
\(\langle(12)\rangle\), and conjugation by \((12)\) inverts \((123)\).  Hence
\[S_3\cong C_3\rtimes C_2.\]
Finally the chain \(S_3\trianglerighteq A_3\trianglerighteq1\) shows that
\(S_3\) is solvable.  Its derived subgroup is \(A_3\), so its abelianization
is \(C_2\).  One group has therefore illustrated actions, cosets, normality,
quotients, conjugacy, semidirect products, composition factors, and
solvability.

\subsection*{The Group \(D_4\)}

The square group \(D_4=\langle r,s\mid r^4=s^2=1,\ srs^{-1}=r^{-1}\rangle\)
has eight elements.  Its geometric action on the four vertices is transitive,
and the stabilizer of a vertex has two elements.  This gives \(8=4\cdot2\),
the orbit--stabilizer formula in a visibly geometric form.

The rotation subgroup \(\langle r\rangle\cong C_4\) is normal because
\(s r^i s^{-1}=r^{-i}\) remains a rotation.  The quotient by rotations is
\(C_2\): it distinguishes orientation-preserving symmetries from reflections.
The relation \(srs^{-1}=r^{-1}\) also displays \(D_4\) as
\(C_4\rtimes C_2\).  This is an especially useful contrast with
\(C_4\times C_2\), whose two factors commute; the dihedral relation gives the
nontrivial twisting.

The center of \(D_4\) is \(\{1,r^2\}\).  Since the quotient by this subgroup
is abelian, the chain \(1\leq\langle r^2\rangle\leq D_4\) is a central series.
Thus \(D_4\) is nilpotent.  Comparing it with \(S_3\), which is solvable but
has trivial center and is not nilpotent, makes the hierarchy
\[\text{abelian}\ \Longrightarrow\ \text{nilpotent}\ \Longrightarrow\
\text{solvable}\]
concrete rather than merely formal.

\begin{remark}[A Study Method]
When a new group-theoretic construction is introduced, revisit \(S_3\) and
\(D_4\).  Ask which subgroups are normal, what the natural actions are, what a
quotient remembers, and whether the group decomposes as a direct or semidirect
product.  This practice develops structural intuition more effectively than a
long list of unrelated calculations.
\end{remark}

\section*{Exercises}
\begin{exercise}[Basic]
Show that the kernel of a permutation representation is the intersection of
all point stabilizers.
\end{exercise}
\begin{exercise}[Core]
Let \(G\) act transitively on a set \(X\), and fix \(x\in X\).  Prove that
the stabilizers of the points of \(X\) are conjugate.  Interpret this result
for the left action of \(G\) on the cosets of a subgroup \(H\).
\end{exercise}
\begin{exercise}[Core]
Compute the conjugacy classes, centralizers, and class equation of \(S_3\).
Explain which features of the computation reflect the geometry of a triangle.
\end{exercise}
\begin{exercise}[Core]
Let \(P\) be a Sylow \(p\)-subgroup of a finite group \(G\).  Prove that
\(P\trianglelefteq G\) if and only if it is the unique Sylow \(p\)-subgroup.
Deduce that any Sylow subgroup whose number is forced to be one by the Sylow
divisibility and congruence conditions is normal.
\end{exercise}
\begin{exercise}[Core]
Let \(G\) be a finite \(p\)-group.  Prove that every nontrivial normal
subgroup of \(G\) meets \(Z(G)\) nontrivially.
\end{exercise}
\begin{exercise}[Further]
Show that \(S_3\cong C_3\rtimes C_2\), but is not a direct product of them.
\end{exercise}
\begin{exercise}[Looking Ahead]
Let \(|G|=p^2\), where \(p\) is prime.  Use the center theorem and the
quotient by a central subgroup of order \(p\) to prove that \(G\) is abelian.
Explain why this does not classify such groups.
\end{exercise}
